\chapter{Design decisions}
\label{ch:design}

Our implementation of \ac{SCITE} is a multi-kernel design. We initially tried to implement the algorithm with a single kernel, but we soon encountered feedback problems that we were not able to resolve inside a single kernel. Therefore, we split the parts of the algorithm with problematic feedback into separate kernels, connected via pipes. Figure \ref{fig:design_schematic} gives an overview of the design: The change proposer kernel receives the current state from the IO kernel and samples all random parameters of a chain step. Everything else like finding the deterministic parameters of the move, computing the resulting tree, computing its likelihood, and evaluating it is done by the ``Tree Scorer'' kernel. It contains four internal loops that have to be completed for every tree. Once the likelihood score of the new tree is computed and the new current state is decided, it is fed back to the IO kernel. Additionally, it keeps track of the most-likely state it encounters and writes this state to a global buffer once the execution is finished. The change proposer kernel and the loops inside the tree scorer kernel all operate independently of each other in a macro-pipeline. The IO kernel dispatches and receives states for the pipeline: Chain states from the tree scorer are fed back to the pipeline as often as requested by the user and when these chains are finished, new initial states are read from global memory. These states are precomputed by the host since we did not want to waste chip space on hardware that is rarely used. It is also not necessary to write the intermediate or final chain states back to global memory since the application only evaluates the best states.

\begin{figure}
    \centering
    \begin{tikzpicture}
        \draw[every rectangle node/.style={draw}]
            node (IOKernel) at (-2,0) {IO Kernel}
            node (ChangeProposer) at (6,0) {Change Proposer}
            node (FindParentLoop) at (2,-3) {Find Parent Loop}
            node (TreeUpdateLoop) at (2,-4) {Tree Update Loop}
            node (TreeScoreMapLoop) at (2,-5) {Tree Score Mapping Loop}
            node (TreeScoreReductionLoop) at (2,-6) {Tree Score Reduction Loop}
            node (InitialStates) at (-5.5,0) {Initial States}
            node (BestStates) at (-5.5,-6) {Best State};

        \node[draw,thick,fit=(FindParentLoop) (TreeUpdateLoop) (TreeScoreMapLoop) (TreeScoreReductionLoop),inner sep=.4cm] (TreeScorer) {};
        \node[above] at (TreeScorer.north) {Tree Scorer};

        \draw[-{Stealth[length=.2cm]}] (IOKernel) 
            -- node[above] (IOCPLink) {current tree \& score}
            (ChangeProposer);
        \draw[-{Stealth[length=.2cm]}] (ChangeProposer) 
            |- node[left,text width=2.5cm,pos=0.1] (CPTSLink) {current tree \& score, move metadata}
            (TreeScorer);
        \draw[-{Stealth[length=.2cm]}] (TreeScorer)
            -| node[right,text width=2.5cm,pos=0.9] (TSIOLink) {current tree \& score}
            (IOKernel);
        \draw[-{Stealth[length=.2cm]}] (InitialStates) 
            -- node[above,text width=2.5cm] {}
            (IOKernel);
        \draw[-{Stealth[length=.2cm]}] (TreeScorer.190) -- +(-1,0) -- +(-1,-1) 
            -- node[below,text width=2.5cm] {}
            (BestStates);

        \node[draw,ultra thick,fit=(IOKernel) (IOCPLink) (ChangeProposer) (CPTSLink) (TreeScorer) (TSIOLink),inner sep=.4cm] (FPGA) {};
        \node[above] at (FPGA.north) {FPGA chip design};

        \node[draw,ultra thick,fit=(InitialStates) (BestStates),inner sep=.4cm] (DDR) {};
        \node[above] at (DDR.north) {Global DDR memory};
    \end{tikzpicture}
    \caption{Schematic of the \ac{ffSCITE} chip design.}
    \label{fig:design_schematic}
\end{figure}

In this chapter, we discuss noteworthy contributions and optimizations. First, we describe how we used both logical properties of mutation trees and technical properties of \acp{FPGA} to efficiently use, modify and transfer mutation trees in the design. Secondly, we describe how we have rewritten the likelihood score computation to make use of wide parallelism available on \acp{FPGA}. Thirdly, we describe our difficulties with move proposals, especially with random number generators, and how we have minimized their impact. Lastly, we describe how we have set up the IO and feedback kernel to saturate the computation pipeline.

\section{Tree encoding and operations}

Our first and most impactful contribution is an improvement to the used mutation tree encoding and the operations on this code. The original \ac{SCITE} implementation \cite{tree2016} uses a parent vector as the canonical data structure of a mutation tree:

\begin{definition}[Parent vector, \cite{tree2016}]
    \label{def:parent_vector}
    Let $T = (V, E, r)$ be a mutation tree. The corresponding parent vector is defined as the sequence $(p_v)_{v \in V} \subseteq V$ with
    \begin{align*}
        p_v &:= \begin{cases}
            p_T(v) & v \neq r \\
            r & v = r 
        \end{cases}
    \end{align*}
\end{definition}

Using this encoding has the obvious advantage that most of the tree moves are simple: The updates in the ``prune and reattach'' and ``swap subtrees'' moves of \textsc{ChainStep} (Figure \ref{alg:scite-step}, lines 21, 36, and 37) are mere constant operations. The update operations in the ``swap nodes'' move are more involved since every edge needs to be visited and checked, but it's still in linear runtime. However, sampling moves and computing the likelihood function requires many connectivity queries: Both the ``prune and reattach'' and the ``swap subtrees'' moves may need to sample a target that is not a descendant of the moved node, and the induced mutation matrix (Definition \ref{def:induced_mutmatrix}) is defined by node connectivity. Therefore, a separate data structure is used to answer these queries quickly:

\begin{definition}[Ancestor matrix, \cite{tree2016}]
    \label{def:ancestor_matrix}
    Let $T = (V, E, r)$ be a mutation tree. The corresponding ancestor matrix is defined as the matrix $A \in \{0,1\}^{|V| \times |V|}$ with
    \begin{align*}
        A_{v,w} &:= \begin{cases}
            1 & v \leadsto_T w \\
            0 & \text{else}
        \end{cases}
    \end{align*}
    for all $v, w \in V$.
\end{definition}

\begin{figure}[p]
    \begin{algorithmic}[1]
        \Function{AncestorMatrix}{$(p_v)_{v \in V} \subseteq V, r \in V$} \Comment $r$ is the root of the tree.
            \State $n \leftarrow |V|$
            \State $a \leftarrow 0 \in \{0,1\}^{n \times n}$
            \ForAll{$w \in V$}
                \State $v \leftarrow w$
                \While{$v \neq r$}
                    \State $a_{v, w} \leftarrow 1$
                    \State $v \leftarrow p_v$
                \EndWhile
                \State $a_{r, w} \leftarrow r$
            \EndFor
            \State \Return $a$
        \EndFunction
    \end{algorithmic}
    \caption{Algorithm to construct an ancestor matrix (Definition \ref{def:ancestor_matrix}) from a parent vector (Definition \ref{def:parent_vector}), \cite{tree2016}}
    \label{alg:ancestor_matrix}
\end{figure}

It is obvious that connectivity queries can be executed in constant time using an ancestor matrix. Jahn et al. \cite{tree2016} also give an algorithm that constructs an ancestor matrix from a parent vector; We have listed it in figure \ref{alg:ancestor_matrix}. Intuitively, it walks up from every node to the root and marks all nodes it encounters as ancestors. \ac{SCITE} uses this algorithm twice per chain step, once before sampling a move and once after the move was executed to compute the likelihood function. It is however very hard to implement in hardware: It is hard to predict how often the inner while-loop is executed. For example, it may be executed $|V|$ times for the leaf of a completely degenerated tree, but it may not be executed at all for the root $r$. Therefore, the outer for-loop needs to be executed serially, which severely limits the performance of the design. However, we were able to eliminate the need to construct ancestor matrices on the device. First of all, we were able to show with theorem \ref{theo:am_reverse} that it is possible to find a node's parent using an ancestor matrix. Ancestor matrices can therefore be used as the canonical data structure to encode mutation trees. More importantly, however, we were able to show with theorem \ref{theo:am_update} that every move of the \textsc{ChainStep} algorithm (Figure \ref{alg:scite-step}) can be executed on an ancestor matrix with linear time and space requirements. This leaves us with two linear, perfectly pipelinable loops, a certain improvement on the previous implementation.

\subsection{Reversing the ancestor matrix construction}

\begin{theorem}
    \label{theo:am_reverse}
    Let $T=(V, E, r)$ be a mutation tree, $a \in \{0, 1\}^{|V| \times |V|}$ be the corresponding ancestor matrix, and $(v, w) \in V^2$ be a node pair. The algorithm \textsc{IsParent} in Figure \ref{alg:is_parent} returns True iff $v = p_T(w)$ given the inputs $V$, $a$ and $(v,w)$. If we implement \textsc{IsParent} as a hardware design and unroll to for-loop at line 8ff., it has a runtime in $O(1)$ and a space usage in $O(|V|)$. Finding a node's parent by iterating over all nodes has a runtime in $O(|V|)$, a space usage $O(|V|)$, and is perfectly pipelinable.
\end{theorem}

The algorithm uses Lemma \ref{lem:am_reverse} to check whether the requested edge exists in the tree, with an additional shortcut for the root. Unrolling the for-loop is possible since we assume an upper bound for the number of genes, and apart from that, the runtime and space requirement should be obvious to see.

\begin{lemma}
    \label{lem:am_reverse}
    Let $T = (V, E)$ be a tree and $v, w \in V$. We have:
    \begin{align*}
        (v, w) \in E \Leftrightarrow (\forall x \in V \setminus \{w\}: x \leadsto v \Leftrightarrow x \leadsto w) \wedge (v \leadsto w)
    \end{align*}
\end{lemma}

\begin{proof}

We first show $\Rightarrow$: We obviously have $v \leadsto w$. Let $x \in V \setminus \{v\}$. Then, we have:
\begin{align*}
    x \leadsto w    &\Rightarrow \exists p = (x, \dots, w) \subseteq E \\
                    &\stackrel{(v, w) \in E}{\Rightarrow} p = (x, \dots, v, w) \subseteq E \\
                    &\Rightarrow p' := (x, \dots, v) \subseteq E \\
                    &\Rightarrow x \leadsto v \\
    x \leadsto v    &\Rightarrow \exists p = (x, \dots, v) \subseteq E \\
                    &\stackrel{(v, w) \in E}{\Rightarrow} p = (x, \dots, v, w) \subseteq E \\
                    &\Rightarrow x \leadsto w
\end{align*}
Now, we show $\Leftarrow$: Let's assume for a contradiction that there is a $y \in V \setminus \{v, w\}$ with $v \leadsto y \leadsto w$. Since we have $y \neq w$, $y \leadsto w \Rightarrow y \leadsto v$ follows. However, this means that we have $v \leadsto y \leadsto v$ and that our tree has a circle, which is a contradiction. Therefore, such a $y$ can not exist and we have in fact $p = (v, \dots, w) = (v, w) \Rightarrow (v, w) \in E$.
\end{proof}

\begin{figure}[p]
    \begin{algorithmic}[1]
        \Function{IsParent}{$V$, $a \in \{0,1\}^{|V| \times |V|}$, $(v, w) \in V^2$}
            \If{$a_{v,w} = 0$}
                \State \Return False
            \EndIf
            \If{$v = w$}
                \State \Return $v = r$ \Comment Per convention, the root is the only node that is also its parent.
            \EndIf
            \ForAll{$x \in V \setminus \{w\}$}
                \If{$a_{x,w} \neq a_{x,v}$}
                    \State \Return False
                \EndIf
            \EndFor
            \State \Return True
        \EndFunction
    \end{algorithmic}
    \caption{Algorithm to query whether an edge exists in a tree, using an ancestor matrix}
    \label{alg:is_parent}
\end{figure}
\section{Likelihood computation}
\label{sec:scoring}

The likelihood score of a mutation tree is used on multiple occasions, mainly to judge whether a newly proposed tree is the best and to compute the probability of accepting a proposed tree. Since every current tree has been a proposed tree at some point, it is possible to carry the likelihood score along with the tree and only compute the likelihood score of the proposed tree in the chain step. Both our and the reference implementation do this in a fundamentally similar way described in algorithm \ref{alg:likelihood}: Both implementations iterate over every cell and tree node and identify the mutations the cell would have if it were attached to this node. These mutations yield a likelihood score for the cell-node pair. As shown in lemma \ref{lem:max_attachment}, we have to choose the attachment node for every cell that yields the highest likelihood, which is then called the likelihood of the cell. These likelihoods are then multiplied together to yield the likelihood of the entire tree.

Both implementations try to minimize the number of floating-point operations to save time and preserve precision. Therefore, they do not directly multiply individual likelihood factors together. Instead, they count the occurrences of true and false positives and negatives, raise the probabilities of these events to the counted numbers, and multiply these powers together. \ac{SCITE} even differentiates between a fast and an accurate computation method: In the fast method, the likelihoods of the cells are directly multiplied together after they are decided, and in the accurate method only the event counts are summed up and the final likelihood score is just computed in the end. \ac{SCITE} first runs the fast method on the newly proposed tree and if the result is close to or higher than the best-known likelihood score, the accurate method is used again. We chose to only implement and use the fast method since our performance would not gain from using both methods: We would need to allocate chip space to both methods and would probably need to run both every time to avoid diverging loops. Our quality benchmark assures us that the accuracy of our implementation is still satisfying.

\todo[inline]{Reference quality benchmark}

Running many power operations however costs a lot of time, and in the case of \acp{FPGA} also a lot of space. Therefore, both implementations work with log-likelihoods instead since power and multiply operations translate to the lightweight multiply and add operations for log-likelihoods. Additionally, log-likelihoods are required since real-world likelihoods often have a very small absolute value that can not be represented even with 64-bit floating point numbers. Using the logarithm of the likelihoods instead brings the absolute values into ranges where even 32-bit floating point numbers are precise enough. Lastly, it is also possible to work with log-likelihoods for most of the algorithm since the logarithm preserves order. Exponentiation of the log-likelihoods is therefore only required when computing the acceptance probability.

There are however two big differences between our implementation and the reference: The reference implementation identifies the mutations of a cell-node pair by walking up the tree from the attachment node to the root using a parent vector (definition \ref{def:parent_vector}), just like in the construction function for ancestor matrices (algorithm \ref{alg:ancestor_matrix}). This is not feasible on \acp{FPGA}, as we have already discussed in section \ref{sec:encoding}. Instead, our implementation iterates over all genes and uses an ancestor matrix to check whether this gene is an ancestor of the attachment node and therefore mutated. The tradeoffs of iterating over all nodes/genes instead of walking up a tree are already discussed in subsection \ref{sec:encoding_implementation}; In summary, it improves the structural performance of the design but it does not scale up for smaller inputs. The next difference is in how the resulting values are reduced: The reference implementation directly reduces found values by always keeping the best-known likelihood of a cell-node pair and an accumulator for the final likelihood. We encountered feedback problems with this approach which inhibited loop pipelining. Therefore, we chose a clear map-reduce approach: We first compute and store the likelihood of every cell-node pair in a buffer, max-reduce the node axis and then sum-reduce the cell axis. Every one of these three operations is executed by separate, independent loops, which are perfectly pipelinable and unrolled by at least one dimension. Additionally, we were able to coalesce and partially unroll the two remaining dimensions of the likelihood mapping loop. However, this obviously increases the memory and chip space requirements of our design. Our resulting log-likelihood computation algorithm is listed as algorithm \ref{alg:likelihoodimpl}.

\begin{algorithm}
    \begin{algorithmic}
        \Function{Likelihood}{$C$, $G$, $T = (V, E, r)$, $d \in \{0,1,2\}^{|C| \times |G|}, \{\alpha, \beta\} \subseteq [0,1]$}
            \State $l_T \leftarrow 1$
            \ForAll{$c \in C$}
                \State $l_c \leftarrow 0$
                \ForAll{$v \in V$}
                    \State $l_{c,v} \leftarrow 1$
                    \ForAll{$g \in G$}
                        \State $l_{c,v} \leftarrow l_{c,v} \cdot \lambda(d_{c,g}, g \leadsto_T v)$
                    \EndFor
                    \If{$l_{c,v} \geq l_c$}
                        \State $l_c \leftarrow l_{c,v}$
                    \EndIf
                \EndFor
                \State $l_T \leftarrow l_T \cdot l_c$
            \EndFor
            \State \Return $l_T$
        \EndFunction
    \end{algorithmic}
    \caption{Formal algorithm to compute $\Lambda_d(T)$ (Definition \ref{def:likelihood}, Lemma \ref{lem:max_attachment}) for a given mutation tree $T$, a mutation data matrix $d$ and error probabilities $\alpha$ and $\beta$.}
    \label{alg:likelihood}
\end{algorithm}

\begin{algorithm}
    \begin{algorithmic}
        \Function{LogLikelihood}{$V$, $C$, $G$, $a \in \{0,1\}^{|V| \times |V|}$, $d \in \{0, 1, 2\}^{|C| \times |G|}$, $\{\alpha, \beta\} \subseteq [0,1]$}
            \State $l_\mathrm{pair} \leftarrow 0 \in (-\infty,0]^{|C| \times |V|}$
            \ForAll{$(c,v) \in C \times V$} \Comment Unroll as much as possible
                \State $\mathrm{occurrences} \leftarrow 0 \in \mathbb{N}_0^{3 \times 2}$
                \ForAll{$g \in G$} \Comment Unroll completely
                    \State $\mathrm{posterior} \leftarrow d[c][g]$
                    \State $\mathrm{prior} \leftarrow a[g][v]$
                    \State $\mathrm{occurrences}[\mathrm{posterior}][\mathrm{prior}] \leftarrow \mathrm{occurrences}[\mathrm{posterior}][\mathrm{prior}] + 1$
                \EndFor
                \State $l_\mathrm{pair}[c][v] \leftarrow \log(1-\alpha) \cdot \mathrm{occurrences}[0][0]$
                \State $l_\mathrm{pair}[c][v] \leftarrow l_\mathrm{pair}[c][v] + \log(\alpha) \cdot \mathrm{occurrences}[1][0]$
                \State $l_\mathrm{pair}[c][v] \leftarrow l_\mathrm{pair}[c][v] + \log(\beta) \cdot \mathrm{occurrences}[0][1]$
                \State $l_\mathrm{pair}[c][v] \leftarrow l_\mathrm{pair}[c][v] + \log(1-\beta) \cdot \mathrm{occurrences}[1][1]$
            \EndFor
            \State $l_\mathrm{cell} \leftarrow 0 \in (-\infty,0]^{|C|}$
            \ForAll{$c \in C$}
                \State $l_\mathrm{cell}[c] \leftarrow \max \{l_\mathrm{pair}[c][v] : v \in V\}$ \Comment Unrolled for-loop in maximum operation
            \EndFor
            \State \Return $\sum_{c \in C} l_\mathrm{cell}$ \Comment Unrolled for-loop in sum
        \EndFunction
    \end{algorithmic}
    \caption{Our \ac{FPGA}-optimized version of algorithm \ref{alg:likelihood} to compute $\log(\Lambda_d(T))$ for a given mutation tree $T$ encoded as an ancestor matrix $a$, a mutation data matrix $d$ and error probabilities $\alpha$ and $\beta$.}
    \label{alg:likelihoodimpl}
\end{algorithm}
\documentclass{cukproposal} % extend to [german] if required

%%%
% Enter Title & Author
\def\Title{Implementing Single Cell Inference of Tumor Evolution of FPGAs}
\def\Author{Jan-Oliver Opdenhövel}
\def\Degree{Bachelor} %or Master
%%%

%%%
% Packages

% Graphical Libs
\usepackage{graphicx}
\usepackage{amsmath, amssymb, amsfonts}% mathematical symbols and the like
\usepackage{amsthm}% definitions, theorems, etc.
\usepackage[colorinlistoftodos]{todonotes}% marking open todos in text/on margins
\usepackage{subfig}% multi-part figures with separate captions per part
\usepackage{url}% render URLs correctly and make them clickable through the hyperref package
\usepackage{longtable}% tables that span multiple pages
\usepackage{booktabs}% tables that actually look good
\usepackage[nolist]{acronym}% consistently use acronyms
\usepackage{pgfgantt} %Ganttchart

%%%
% Defs

% Math Defs
\renewcommand{\epsilon}{\varepsilon}
\DeclareMathOperator*{\argmax}{arg\,max}
\DeclareMathOperator*{\argmin}{arg\,min}

\newcommand*{\bfrac}[2]{\genfrac{\lbrace}{\rbrace}{0pt}{}{#1}{#2}} %Fraction withour bar

\newacro{template}[UPB-CS-TT]{Paderborn University Computer Science thesis template}
\newacro{SCITE}[SCITE]{``Single Cell Inference of Tumor Evolution''}
\newacro{MCMC}[MCMC]{Monte Carlo Markov Chain}
\newacro{FPGA}[FPGA]{Field-Programmable Gate Array}
\newacro{HPC}[HPC]{High-Performance Computing}
\newacro{RNG}[RNG]{random number generator}

\DeclareMathOperator{\testop}{top}

% Customizing GanttChart
% Taken from https://git.cs.uni-paderborn.de/syssec/templates/exposee
\definecolor{barblue}{RGB}{153,204,254}
\definecolor{groupblue}{RGB}{51,102,254}
\definecolor{linkred}{RGB}{165,0,33}
\renewcommand\sfdefault{phv}
\renewcommand\mddefault{mc}
\renewcommand\bfdefault{bc}
\setganttlinklabel{s-s}{START-TO-START}
\setganttlinklabel{f-s}{FINISH-TO-START}
\setganttlinklabel{f-f}{FINISH-TO-FINISH}
\sffamily
\ganttset
{%
	y unit chart=0.6cm,
	canvas/.append style={fill=none, draw=black!35, line width=.5pt},
	hgrid={*1{draw=black!35, line width=.1pt}},
	hgrid style/.style={draw=black!35,line width=.1pt},
	vgrid={*1{draw=black!35, line width=.1pt}},
	today rule/.style={draw=black!64,dash pattern=on 3.5pt off 4.5pt,line width=1.5pt},
	today label font=\small\bfseries,
	title/.style={draw=none, fill=none},
	title label font=\bfseries\footnotesize,
	title label node/.append style={below=7pt},
	progress label text={\pgfmathprintnumber[precision=0, verbatim]{#1}\%},
	include title in canvas=false,
	bar label font=\mdseries\small\color{black!70},
	bar label node/.append style={left=0cm},
	bar/.append style={draw=none, fill=black!63},
	bar incomplete/.append style={fill=barblue},
	bar progress label font=\mdseries\footnotesize\color{black!70},
	bar height=.8,
	bar top shift=0.1,
	group top shift=0.1,
	group incomplete/.append style={fill=groupblue},
	group left shift=0,
	group right shift=0,
	group height=.8,
	group peaks tip position=0,
	group label node/.append style={left=.6cm},
	group progress label font=\bfseries\small,
	link/.style={-latex, line width=1.5pt, linkred},
	link label font=\scriptsize\bfseries,
	link label node/.append style={below left=-2pt and 0pt}
}

\begin{document}

\bibliographystyle{alpha}

\thispagestyle{plain}

\includegraphics[width=7cm]{figures/upb_logo}

\bigskip

\hrule width \hsize height 2pt \kern 1mm \hrule width \hsize
\vspace*{6pt}

\begin{center}
	\Large\textbf{\Title}
\end{center}

\noindent
\begin{minipage}{6cm}
	\large{
		\ifgerman
			Proposal zur \Degree arbeit
		\else
			\Degree 's Thesis Proposal
		\fi
	}

	\ifgerman
		\large von
	\else
		\large by
	\fi
	\Author
\end{minipage}
\hfill
\begin{minipage}{6cm}
	\flushright
	\normalsize Research Group

	High-Performance Computing
\end{minipage}

\bigskip

\hrule width \hsize \kern 1mm \hrule width \hsize height 2pt


\chapter{Introduction}
\label{ch:introduction}

\section{Motivation and Goal of the Thesis}

Cancer is a widespread and often lethal disease\cite{10.1001/jamaoncol.2021.6987} where body cells mutate in a way that increases their cell division speed as well as their lifetime while also evading natural cancer prevention mechanisms. There are many possible treatment methods like surgery, chemotherapy, or radiation therapy, but their effectiveness often depends on the exact type of tumor. The complexity of choosing the correct therapy is increased since the tumor cells may mutate again and form metastases in different parts of the body. These developments however can be anticipated by analysing the genome of a tumor and finding new, emerging subclones that may require different treatments even before they are dominant.

Bulk sequencing of tumor cells and analysing the found mutations is already widely used, but often misses out on smaller, upcoming variants since they are averaged out in the mass of cells. Therefore, research has been done to utilise single cell sequencing. With this technique, the exact genome of individual cells can be identified and compared against other cells. However, it also comes with high error rates and parts of a cell's genome are often lost during the process. It is therefore impossible to compute an exact phylogenetic tree with this data. Instead, one has to resort to find the most plausible one.

One algorithm to compute the most likely phylogeny from a set of erroneous cell genomes is \ac{SCITE}\cite{tree2016}. The original authors provided a functional, but rather unoptimized implementation of their algorithm. While I know of an optimized but unpublished implementation of \ac{SCITE} for CPUs, there has not been an efficient implementation for \acp{FPGA} to my knowledge. \acp{FPGA} are computer chips that contain a lattice of logic, computation, and memory units that are connected via programmable connections. They basically emulate a complete computer chip and have already been used extensively for chip prototyping and verification. In recent years they have also become interesting as computation accelerators for \ac{HPC} clusters due to their low power consumption compared to their performance. However, developing efficient \ac{FPGA} designs is tedious since compiling them often takes multiple hours and naive designs are often multiple orders of magnitudes worse than optimized ones. \todo[inline]{Add notes about usage}

However, with the introduction of Intel's oneAPI, it has become possible to develop \ac{FPGA} designs using C++ templates, a form of polymorphism that is very suitable for \ac{FPGA} designs: Datatypes and functions are implemented with placeholders that can be replaced with actual types and constants during compilation. When set up correctly, one can therefore implement different solutions to the same problem, integrate the different solutions in the build system, compile the different versions in parallel and evaluate the resulting designs side by side. I have already employed this method while developing StencilStream, a Generic Stencil Simulation Library for \acp{FPGA}, and it significantly improved my productivity.

My goal for this thesis is therefore to develop an efficient implementation of SCITE for \acp{FPGA} using oneAPI and the aforementioned development pattern. This implementation should be at least as fast as the original author's implementation, but possibly as fast as the optimized CPU version mentioned previously or even faster. With this, I want to showcase the increased productivity of \ac{FPGA} development with oneAPI, that even an undergraduate student can develop efficient FPGA designs in the limited amount of time of a Bachelor's Thesis.

\section{Problem Formulation}

The input is a function $D: C \times G \rightarrow \{0, 1, ?\}$, where $C$ is the set of sampled cells and $G$ is a set of genome positions. If for a cell $c \in C$ and a genome position $g \in G$ we have $D(c, g) = 1$, the genome position $g$ has been observed as mutated in cell $c$, and if we have $D(c, g) = 0$, the genome position $g$ has been observed in it's common form. There is also a third option $D(c, g) = ?$, which means that the genome position has not been found at all and no statement could be made whether it is mutated or not. As the existance of the $?$ option may indicate, this function is noisy and may contain errors. It is assumed that errors occur independently from each other and if we let $E: C \times G \rightarrow \{0, 1\}$ be the true state of the cells, we have
\begin{align*}
    \mathbb{P}(D(c, g) = 0 \mid E(c, g) = 0) &= 1 - \alpha & \mathbb{P}(D(c, g) = 1 \mid E(c, g) = 0) &= \alpha \\
    \mathbb{P}(D(c, g) = 0 \mid E(c, g) = 1) &= \beta & \mathbb{P}(D(c, g) = 1 \mid E(c, g) = 1) &= 1 - \beta
\end{align*}
for some error rates $\alpha, \beta \in (0, 1)$. Error modeling for $D(c, g) = ?$ is not given by the authors. Missing data is always ignored in probability calculations.

The goal now is to find a mutation tree, a cell attachment and error rates with a high likelihood of being correct. A mutation tree is a rooted tree $T = (G \cup \{r\}, E, r)$ where $r$ is the root of the tree and the other nodes are genome positions. It expresses that the mutation of $g \in G$ has occured after all it's ancestors, prior to it's successors and independently from all unrelated genome positions. A cell attachment is a simple function $\sigma: C \rightarrow G$ and expresses that a cell $c \in C$ has the mutation status just after $\sigma(c)$ was mutated. Therefore, a cell $c$ has mutations at all genome positions that are ancestors of $\sigma(c)$, but at no other positions. This model is sensible since the authors of \ac{SCITE} make the ``infinite sites'' assumption, which states that a genome position mutates exactly once in the complete history of the tumor and does not mutate back.

With the mutation tree $T$ and a cell attachment $\sigma$, we can now build a function $e_{(T, \sigma)}: C \times G \rightarrow \{0, 1\}$ that expresses the made statements of the cell's true mutation status. Formally, we set
\begin{align*}
    e_{(T, \sigma)}(c, g) &:= \begin{cases}
        1 & g \in \mathrm{Anc}_T(\sigma(c)) \\
        0 & \text{else} \\
    \end{cases}
\end{align*}
where $\mathrm{Anc}_T(g)$ is the set of ancestors of a genome position $g \in G$ with respect to $T$. Therefore, the likelihood of a tree and cell attachment given the observed data $d$ and the error rates $\alpha$ and $\beta$ is
\begin{align*}
    \prod_{c \in C}\prod_{g \in G} \mathbb{P}_{(\alpha, \beta)}(D(c, g) = d(c, g) \mid E(c, g) = e_{(T, \sigma)}(c, g))
\end{align*}
Since pairs of $(c, g) \in C \times G$ with $d(c, g) = ?$ are ignored, the probabilities in the product above are always set to 1. As stated before, the goal is now to find the tuple $(T, \sigma, \alpha, \beta)$ with the highest likelihood.

\section{Algorithm introduction}

\ac{SCITE} is a \ac{MCMC} algorithm. A Monte Carlo algorithm runs a random experiment that produces possible solutions to a problem and evaluates how well the generated solutions solve the problem. This operation of generating and evaluating a solution is then repeated multiple times and the result is the best solution the algorithm has encountered. In theory, it would suffice if the experiment had a probability greater than zero to produce a good solution, but in order to improve the solution quality and to reduce the required repetitions, one would use an experiment that produces the best solutions with a higher probability than worse solutions and that can be repeated quickly. As the name implies, \ac{MCMC} algorithms simulate a Markov Chain to produce solutions. The advantage of using Markov Chains is that the next sample may depend on the previous one. The algorithm may therefore introduce small changes to the solution. This is obviously faster than generating a new solution from scratch and if the current sample is already a good solution, the change may preserve some of its quality. However, a designer of a \ac{MCMC} algorithm has to make sure that the chain actually converges on the desired distribution.

\section{Previous Work}

\section{Work Schedule}

\vspace*{.5cm}
\ifgerman
Paderborn, den
\else
Date:
\fi
\vspace*{1cm}
\begin{center}
	\begin{tabular}{l p{0.1\textwidth} r}
	  \cline{1-1} \cline{3-3}
	  \begin{minipage}[t]{0.4\textwidth}
	    \centering
	    Dr. Tobias Kenter
		\end{minipage}
		&
		\begin{minipage}[t]{0.2\textwidth}
		\end{minipage}
		&
		\begin{minipage}[t]{0.4\textwidth}
		  \centering
		  \Author
		\end{minipage}
	\end{tabular}
\end{center}


%%
\bibliography{references}

\end{document}

\section{Input, Output and Feedback}
\label{sec:io}

The change proposer kernel and the internal loops of the tree scorer kernel form a macro-pipeline that processes states in independent steps: For example, the change proposer may already sample a move for a new state while the tree update loop computes the new ancestor matrix of the previous state. This macro-pipeline has to be filled at all times to utilize the entire design and achieve maximum performance. Ensuring this is the task of the IO kernel.

From an application's point of view, the design is supposed to iterate $n_\mathrm{steps}$ times on the same state and then generate a new, completely random initial state to work with, up to $n_\mathrm{chains}$ times. Early on, we decided to not generate the initial states on the device since this is a rare task compared to the chain steps. We, therefore, assumed that allocating precious resources to this task would not accelerate the application. Instead, our application's host generates $n_\mathrm{chains}$ initial states and stores them in global buffers that are transferred to the \ac{FPGA}'s memory. A na\"ive approach would load a state from the global buffer, send it to the change proposer kernel and wait until the tree scorer is finished to write the result back. The disadvantages of this approach are obvious and lead to low hardware utilization and high latency due to the external memory accesses. However, we were able to exploit three properties of the problem and the hardware design to optimize the occupancy: First, the intermediate and final states of the chains are irrelevant to the algorithm. Only the current state is needed to compute the next state and only the best state needs to be saved and exported. Once a current state has been replaced, it is not needed anymore. Second, nothing in the macro-pipeline needs to know to which step of which chain the processed state belongs. It simply executes the same operations on the initial state of the first chain as it does on the 42,000th step of the 17th chain. Third, the macro-pipeline has a capacity. This means that the IO kernel can send a certain number of states to the change proposer before it needs to start to receive states from the tree scorer. From the IO kernel's point of view, these states are ``stored'' in the pipeline.

Our final IO kernel (algorithm \ref{alg:io-kernel}) works as follows: First, it assumes that the pipeline has a capacity of $c$ states. Therefore, it initially sends $c$ states from the initial states buffer to the change proposer and waits for the tree scorer to finish the first step. Then, the resulting state from the tree scorer is fed back to the change proposer as often as requested by the user. Once the requested steps are completed, the final chain state is discarded and a new initial state is sent to the change proposer instead. Once all steps of all chains are completed, the kernel flushes the pipeline and halts. This design has multiple advantages over the na\"ive approach: Firstly, this design is simple. Its logic is simple and its actions are easy to predict, and there are no loop-carried dependencies. Its biggest advantage however is that it only reads from global memory and does not write back. This simplifies linking as well as routing and placing, which the current \ac{FPGA} compilers still have problems with sometimes and takes a lot of time. This has been entirely avoided here. Lastly, it might also reduce the energy consumption of the design since data only moves within the chip, but we have not analyzed this. However, this design has the disadvantage that one needs to compile the design multiple times to find the right capacity $c$, which takes additional effort and \ac{CPU}-hours for compilation.

\begin{algorithm}
    \begin{algorithmic}
        \ForAll{$i \leq n_\mathrm{chains} \cdot n_\mathrm{steps} + c$}
            \If{$i \geq c$}
                \State $s_\mathrm{out} \leftarrow \text{Read state from tree scorer kernel}$
            \EndIf
            \If{$i < n_\mathrm{chains} \cdot n_\mathrm{steps}$}
                \If{$i \mod (n_\mathrm{steps} \cdot c) < c$}
                    \State $s_\mathrm{in} \leftarrow \text{Read next initial state from global memory}$
                \Else
                    \State $s_\mathrm{in} \leftarrow s_\mathrm{out}$
                \EndIf
                \State Send $s_\mathrm{in}$ to change proposer kernel
            \EndIf
        \EndFor
    \end{algorithmic}
    \caption{Behavioral code of the IO kernel, assuming that the pipeline has a capacity of $c$ states and that the user has requested to simulate $n_\mathrm{chains}$ chains with $n_\mathrm{steps}$ steps each.}
    \label{alg:io-kernel}
\end{algorithm}